%!TeX root=../pridetop.tex

\chapter[Chapter \thechapter]{}
	
	\begin{figure}[t!]
\centering
\includegraphics[width=.7\linewidth]{21top}
\captionlistentry{Headpiece to Chapter \thechapter}
\end{figure}

\lettrine[lines=6,image=true]{initials/chap21t}{he}  discussion of Mr Collins's offer was now nearly at an end, and Elizabeth had only to suffer from the uncomfortable feelings necessarily attending it, and occasionally from some peevish allusion of her mother. As for the gentleman himself, \textit{his} feelings were chiefly expressed, not by embarrassment or dejection, or by trying to avoid her, but by stiffness of manner and resentful silence. He scarcely ever spoke to her; and the assiduous attentions which he had been so sensible of himself were transferred for the rest of the day to Miss Lucas, whose civility in listening to him was a seasonable relief to them all, and especially to her friend.

The morrow produced no abatement of Mrs Bennet's ill humour or ill health. Mr Collins was also in the same state of angry pride. Elizabeth had hoped that his resentment might shorten his visit, but his plan did not appear in the least affected by it. He was always to have gone on Saturday, and to Saturday he still meant to stay.

After breakfast, the girls walked to Meryton, to inquire if Mr Wickham were returned, and to lament over his absence from the Netherfield ball. He joined them on their entering the town, and attended them to their aunt's, where his regret and vexation and the concern of everybody were well talked over. To Elizabeth, however, he voluntarily acknowledged that the necessity of his absence \textit{had} been self-imposed.

»I found,« said he, »as the time drew near, that I had better not meet Mr Darcy;—that to be in the same room, the same party with him for so many hours together, might be more than I could bear, and that scenes might arise unpleasant to more than myself.«

She highly approved his forbearance; and they had leisure for a full discussion of it, and for all the commendations which they civilly bestowed on each other, as Wickham and another officer walked back with them to Longbourn, and during the walk he particularly attended to her. His accompanying them was a double advantage: she felt all the compliment it offered to herself; and it was most acceptable as an occasion of introducing him to her father and mother.

Soon after their return, a letter was delivered to Miss Bennet; it came from Netherfield, and was opened immediately. The envelope contained a sheet of elegant, little, hot-pressed paper, well covered with a lady's fair, flowing hand; and Elizabeth saw her sister's countenance change as she read it, and saw her dwelling intently on some particular passages. Jane recollected herself soon; and putting the letter away, tried to join, with her usual cheerfulness, in the general conversation: but Elizabeth felt an anxiety on the subject which drew off her attention even from Wickham; and no sooner had he and his companion taken leave, than a glance from Jane invited her to follow her upstairs. When they had gained their own room, Jane, taking out her letter, said, »This is from Caroline Bingley: what it contains has surprised me a good deal. The whole party have left Netherfield by this time, and are on their way to town; and without any intention of coming back again. You shall hear what she says.«

She then read the first sentence aloud, which comprised the information of their having just resolved to follow their brother to town directly, and of their meaning to dine that day in Grosvenor Street, where Mr Hurst had a house. The next was in these words:—»I do not pretend to regret anything I shall leave in Hertfordshire except your society, my dearest friend; but we will hope, at some future period, to enjoy many returns of that delightful intercourse we have known, and in the meanwhile may lessen the pain of separation by a very frequent and most unreserved correspondence. I depend on you for that.« To these high-flown expressions Elizabeth listened with all the insensibility of distrust; and though the suddenness of their removal surprised her, she saw nothing in it really to lament: it was not to be supposed that their absence from Netherfield would prevent Mr Bingley's being there; and as to the loss of their society, she was persuaded that Jane must soon cease to regard it in the enjoyment of his.

»It is unlucky,« said she, after a short pause, »that you should not be able to see your friends before they leave the country. But may we not hope that the period of future happiness, to which Miss Bingley looks forward, may arrive earlier than she is aware, and that the delightful intercourse you have known as friends will be renewed with yet greater satisfaction as sisters? Mr Bingley will not be detained in London by them.«

Caroline decidedly says that none of the party will return into Hertfordshire this winter. I will read it to you.

»When my brother left us yesterday, he imagined that the business which took him to London might be concluded in three or four days; but as we are certain it cannot be so, and at the same time convinced that when Charles gets to town he will be in no hurry to leave it again, we have determined on following him thither, that he may not be obliged to spend his vacant hours in a comfortless hotel. Many of my acquaintance are already there for the winter: I wish I could hear that you, my dearest friend, had any intention of making one in the crowd, but of that I despair. I sincerely hope your Christmas in Hertfordshire may abound in the gaieties which that season generally brings, and that your beaux will be so numerous as to prevent your feeling the loss of the three of whom we shall deprive you.«

»It is evident by this,« added Jane, »that he comes back no more this winter.«

»It is only evident that Miss Bingley does not mean he \textit{should}.«

»Why will you think so? It must be his own doing; he is his own master. But you do not know \textit{all}. I \textit{will} read you the passage which particularly hurts me. I will have no reserves from \textit{you}. »Mr Darcy is impatient to see his sister; and to confess the truth, \textit{we} are scarcely less eager to meet her again. I really do not think Georgiana Darcy has her equal for beauty, elegance, and accomplishments; and the affection she inspires in Louisa and myself is heightened into something still more interesting from the hope we dare to entertain of her being hereafter our sister. I do not know whether I ever before mentioned to you my feelings on this subject, but I will not leave the country without confiding them, and I trust you will not esteem them unreasonable. My brother admires her greatly already; he will have frequent opportunity now of seeing her on the most intimate footing; her relations all wish the connection as much as his own; and a sister's partiality is not misleading me, I think, when I call Charles most capable of engaging any woman's heart. With all these circumstances to favour an attachment, and nothing to prevent it, am I wrong, my dearest Jane, in indulging the hope of an event which will secure the happiness of so many?« What think you of \textit{this} sentence, my dear Lizzy?« said Jane, as she finished it. »Is it not clear enough? Does it not expressly declare that Caroline neither expects nor wishes me to be her sister; that she is perfectly convinced of her brother's indifference; and that if she suspects the nature of my feelings for him she means (most kindly!) to put me on my guard. Can there be any other opinion on the subject?«

»Yes, there can; for mine is totally different. Will you hear it?«

»Most willingly.«

»You shall have it in a few words. Miss Bingley sees that her brother is in love with you and wants him to marry Miss Darcy. She follows him to town in the hope of keeping him there, and tries to persuade you that he does not care about you.«

Jane shook her head.

»Indeed, Jane, you ought to believe me. No one who has ever seen you together can doubt his affection; Miss Bingley, I am sure, cannot: she is not such a simpleton. Could she have seen half as much love in Mr Darcy for herself, she would have ordered her wedding clothes. But the case is this:—we are not rich enough or grand enough for them; and she is the more anxious to get Miss Darcy for her brother, from the notion that when there has been \textit{one} inter-marriage, she may have less trouble in achieving a second; in which there is certainly some ingenuity, and I dare say it would succeed if Miss de Bourgh were out of the way. But, my dearest Jane, you cannot seriously imagine that, because Miss Bingley tells you her brother greatly admires Miss Darcy, he is in the smallest degree less sensible of \textit{your} merit than when he took leave of you on Tuesday; or that it will be in her power to persuade him that, instead of being in love with you, he is very much in love with her friend.«

»If we thought alike of Miss Bingley,« replied Jane, »your representation of all this might make me quite easy. But I know the foundation is unjust. Caroline is incapable of wilfully deceiving anyone; and all that I can hope in this case is, that she is deceived herself.«

»That is right. You could not have started a more happy idea, since you will not take comfort in mine: believe her to be deceived, by all means. You have now done your duty by her, and must fret no longer.«

»But, my dear sister, can I be happy, even supposing the best, in accepting a man whose sisters and friends are all wishing him to marry elsewhere?«

»You must decide for yourself,« said Elizabeth; »and if, upon mature deliberation, you find that the misery of disobliging his two sisters is more than equivalent to the happiness of being his wife, I advise you, by all means, to refuse him.«

»How can you talk so?« said Jane, faintly smiling; »you must know, that, though I should be exceedingly grieved at their disapprobation, I could not hesitate.«

»I did not think you would; and that being the case, I cannot consider your situation with much compassion.«

»But if he returns no more this winter, my choice will never be required. A thousand things may arise in six months.«

The idea of his returning no more Elizabeth treated with the utmost contempt. It appeared to her merely the suggestion of Caroline's interested wishes; and she could not for a moment suppose that those wishes, however openly or artfully spoken, could influence a young man so totally independent of everyone.

She represented to her sister, as forcibly as possible, what she felt on the subject, and had soon the pleasure of seeing its happy effect. Jane's temper was not desponding; and she was gradually led to hope, though the diffidence of affection sometimes overcame the hope, that Bingley would return to Netherfield, and answer every wish of her heart.

They agreed that Mrs Bennet should only hear of the departure of the family, without being alarmed on the score of the gentleman's conduct; but even this partial communication gave her a great deal of concern, and she bewailed it as exceedingly unlucky that the ladies should happen to go away just as they were all getting so intimate together. After lamenting it, however, at some length, she had the consolation of thinking that Mr Bingley would be soon down again, and soon dining at Longbourn; and the conclusion of all was the comfortable declaration, that, though he had been invited only to a family dinner, she would take care to have two full courses.
